\section{The problem}
\label{sec:problem}

Ask whether the grammaticality judgments in a syntax paper needed consent and ethics review. The answer turns on what the question hides: whether the study is about that person, not just whether a person produced the judgment. Expert grammaticality judgments occupy an unstable methodological category. In one description, they are human responses to linguistic stimuli, and so look like participant data. In another, they are scholarly evaluations of linguistic materials, and so look like expert ratings, editorial judgments, coding decisions, or psychometric scoring.

The pull toward the first description has a respectable source. A strand of experimental syntax has argued that informal intuitions, often the author's own, are a weak evidential base, and that acceptability should be measured on samples of naive speakers under controlled conditions \citep{gibsonfedorenko2010,gibsonfedorenko2013,gibson2011mturk}. On that view the judges are a sample, their responses are the data, and the design is participant research. The worry is real, and the corrective has improved the field.

The corrective answers a measurement question, not a role question. That informal intuitions can be unreliable is a reason to calibrate, sample, or formalise the measurement, and the reliability of informal judgments has itself been tested and largely upheld \citep{sprouse2013,schutze2016}. It is not a reason to treat everyone who supplies a grammaticality judgment as the object of study. Reliability is a property of the instrument; it does not fix what the study is about.

So the answer is not that expert judgment escapes scrutiny, but that the scrutiny should follow the role. A judge is a participant when the question concerns that person's reaction, processing, background, identity, or distribution of responses. A judge is an evaluator when the question concerns the linguistic status of an item and expertise is used to assess it against disciplinary standards. Misclassifying costs both ways: treating evaluation as sampling overstates what the judgment shows, and treating sampling as evaluation strips persons of protections they are owed.

The target claim is narrow. Expert judges of grammaticality should not be treated as participants merely because their judgments are human acts. The first question is what role the judgment plays in the evidential design.
